\documentclass[UTF8]{ctexart}
\usepackage[a4paper,margin=2.5cm]{geometry}
\usepackage{booktabs}
\usepackage{longtable}
\usepackage{array}
\usepackage{hyperref}
\usepackage{xcolor}

\newcolumntype{L}[1]{>{\raggedright\arraybackslash}p{#1}}
\emergencystretch=3em

\hypersetup{
  colorlinks=true,
  linkcolor=blue,
  urlcolor=blue
}

\title{TE-KG 数据库现状说明}
\author{TE-KG Project}
\date{2026-05-17}

\begin{document}
\maketitle

\section{范围}

本文记录当前 TE-KG 项目的数据库与数据资产状态，覆盖 Neo4j、MySQL、主要运行时数据文件、非智能问答 API 和构建脚本。本文不覆盖智能问答、agent 编排、LLM 调用链路和 agent 插件内部实现。

\section{总体结构}

当前项目是本地 PHP、JavaScript、Neo4j、MySQL 混合应用。运行时页面仍位于项目根目录，数据资产主要位于 \verb|data| 目录，构建和检查脚本位于 \verb|scripts| 目录。

\begin{itemize}
  \item Neo4j 当前运行目标为 \texttt{tekg3}。
  \item MySQL expression 数据库为 \texttt{tekg\_expression}。
  \item PHP 路径入口为 \verb|path_config.php|。
  \item PHP 运行时 Neo4j 配置入口为 \verb|api/runtime_config.php|。
  \item 浏览器端路径入口为 \verb|assets/js/tekg_paths.php|。
  \item Python 路径入口为 \verb|scripts/path_helpers.py|。
\end{itemize}

\section{Neo4j 数据库}

\subsection{连接目标}

当前本地配置位于 \verb|api/config.local.php|，Neo4j HTTP endpoint 指向：

\begin{verbatim}
http://127.0.0.1:7474/db/tekg3/tx/commit
\end{verbatim}

运行时代码不再在各 API 文件中各自维护 \texttt{tekg2} 或 \texttt{tekg21} 默认值。缺少本地配置或环境变量时，应由 \verb|api/runtime_config.php| 抛出明确配置错误。

\subsection{节点标签计数}

\begin{longtable}{lr}
\toprule
Label & Count \\
\midrule
\endhead
Function & 3683 \\
Paper & 2308 \\
Gene & 1280 \\
Protein & 1089 \\
DiseaseCategory & 767 \\
Disease & 676 \\
RNA & 588 \\
Mutation & 377 \\
Pharmaceutical & 293 \\
TE & 225 \\
Toxin & 67 \\
Lipid & 26 \\
Peptide & 23 \\
Carbohydrate & 12 \\
NonHumanTE & 1 \\
\bottomrule
\end{longtable}

\subsection{关系类型计数}

\begin{longtable}{lr}
\toprule
Relationship type & Count \\
\midrule
\endhead
BIO\_RELATION & 12444 \\
HAS\_SUBCATEGORY & 744 \\
CLASSIFIED\_AS & 436 \\
SUBFAMILY\_OF & 72 \\
\bottomrule
\end{longtable}

\subsection{主要关系形态}

当前 \texttt{BIO\_RELATION} 是主要生物关系类型，具体语义放在关系属性 \texttt{predicate} 中。计数较高的关系形态包括：

\begin{longtable}{lllr}
\toprule
Type & Source label & Target label & Count \\
\midrule
\endhead
BIO\_RELATION & TE & Function & 2841 \\
BIO\_RELATION & Protein & Function & 1377 \\
BIO\_RELATION & Function & Function & 1194 \\
BIO\_RELATION & Function & Disease & 785 \\
HAS\_SUBCATEGORY & DiseaseCategory & DiseaseCategory & 744 \\
BIO\_RELATION & TE & Disease & 526 \\
CLASSIFIED\_AS & Disease & DiseaseCategory & 436 \\
BIO\_RELATION & TE & Gene & 414 \\
BIO\_RELATION & Protein & TE & 407 \\
BIO\_RELATION & Protein & Protein & 406 \\
SUBFAMILY\_OF & TE & TE & 72 \\
\bottomrule
\end{longtable}

\subsection{核心节点属性}

\begin{longtable}{L{3.4cm}L{10.8cm}}
\toprule
Label & Observed property keys \\
\midrule
\endhead
TE & \texttt{name}, \texttt{description}, \texttt{source\_group}, \texttt{taxonomy\_group}, \texttt{taxonomy\_status}, \texttt{taxonomy\_source}, \texttt{taxonomy\_canonical\_name}, \texttt{taxonomy\_class}, \texttt{taxonomy\_subclass}, \texttt{taxonomy\_order}, \texttt{taxonomy\_superfamily}, \texttt{taxonomy\_family}, \texttt{taxonomy\_subclade}, \texttt{is\_leaf\_standard}, \texttt{homepage\_chart\_included} \\
Paper & \texttt{name}, \texttt{description}, \texttt{source\_group}, \texttt{pmid} \\
Disease & \texttt{name}, \texttt{description}, \texttt{source\_group}, \texttt{disease\_class} \\
DiseaseCategory & \texttt{name}, \texttt{source\_group}, \texttt{category\_level}, \texttt{category\_node\_id}, \texttt{category\_label}, \texttt{top\_category}, \texttt{is\_leaf}, \texttt{path\_key} \\
Gene & \texttt{name}, \texttt{description}, \texttt{source\_group}, \texttt{reclassified\_reason} \\
Other biochemical labels & \texttt{name}, \texttt{description}, \texttt{source\_group}. This includes Protein, RNA, Function, Mutation, Pharmaceutical, Toxin, Lipid, Peptide, and Carbohydrate. \\
\bottomrule
\end{longtable}

\subsection{关系属性}

\begin{itemize}
  \item \texttt{BIO\_RELATION}: commonly has \texttt{predicate}, \texttt{pmids}, \texttt{description}, and \texttt{source\_group}.
  \item \texttt{SUBFAMILY\_OF}: may have \texttt{child\_tree\_label}, \texttt{parent\_tree\_label}, \texttt{tree\_reference}, and \texttt{source}.
  \item \texttt{CLASSIFIED\_AS}: has disease classification metadata such as \texttt{source\_disease\_name}, \texttt{mapping\_type}, \texttt{top\_category}, \texttt{leaf\_category\_label}, and \texttt{path\_depth}.
  \item \texttt{HAS\_SUBCATEGORY}: currently has \texttt{source\_group}.
\end{itemize}

\subsection{约束与索引}

当前数据库有节点唯一性约束，主要覆盖核心实体的名称字段：

\begin{longtable}{lll}
\toprule
Constraint & Label & Property \\
\midrule
\endhead
te\_name\_unique & TE & name \\
function\_name\_unique & Function & name \\
gene\_name\_unique & Gene & name \\
protein\_name\_unique & Protein & name \\
rna\_name\_unique & RNA & name \\
disease\_name\_unique & Disease & name \\
mutation\_name\_unique & Mutation & name \\
pharmaceutical\_name\_unique & Pharmaceutical & name \\
toxin\_name\_unique & Toxin & name \\
lipid\_name\_unique & Lipid & name \\
peptide\_name\_unique & Peptide & name \\
carbohydrate\_name\_unique & Carbohydrate & name \\
paper\_pmid\_unique & Paper & pmid \\
disease\_category\_id\_unique & DiseaseCategory & category\_node\_id \\
\bottomrule
\end{longtable}

\section{TE Taxonomy 状态}

\subsection{权威来源}

当前运行时 TE taxonomy 权威来源是 Neo4j \texttt{tekg3} 中 \texttt{TE} 节点的 \texttt{taxonomy\_*} 属性。旧的 lineage JSON/CSV 和旧 G6 demo data 不再作为运行时真相源。

\subsection{分类计数}

\begin{longtable}{lr}
\toprule
Metric & Count \\
\midrule
\endhead
TE nodes & 225 \\
TE nodes with taxonomy\_group & 225 \\
TE nodes with taxonomy\_class & 192 \\
Homepage chart nodes & 154 \\
Leaf standard nodes & 188 \\
\bottomrule
\end{longtable}

\begin{longtable}{lr}
\toprule
taxonomy\_group & Count \\
\midrule
\endhead
standard & 138 \\
A & 36 \\
B & 35 \\
C & 16 \\
\bottomrule
\end{longtable}

\begin{longtable}{lr}
\toprule
taxonomy\_class & Count \\
\midrule
\endhead
Retrotransposons & 156 \\
DNA Transposons & 36 \\
missing & 33 \\
\bottomrule
\end{longtable}

\subsection{主要 order 与 superfamily}

\begin{longtable}{lr}
\toprule
taxonomy\_order & Count \\
\midrule
\endhead
missing & 111 \\
LTR Retrotransposons & 90 \\
Non-LTR Retrotransposons (LINEs) & 24 \\
\bottomrule
\end{longtable}

\begin{longtable}{lr}
\toprule
taxonomy\_superfamily & Count \\
\midrule
\endhead
ERV1 & 41 \\
missing & 41 \\
ERV3 (HERV-R) & 26 \\
SINE1/7SL (Alu) & 23 \\
ERV2 (HERVK) & 21 \\
L1 (LINE-1) & 21 \\
Tc1/Mariner & 13 \\
hAT & 13 \\
SINE-R & 10 \\
SINE2/tRNA & 5 \\
\bottomrule
\end{longtable}

\subsection{标准化报告}

当前标准化报告位于：

\begin{verbatim}
data/processed/tekg3_taxonomy_standardization_report.json
\end{verbatim}

报告现在区分输入记录与合并后节点：

\begin{itemize}
  \item \texttt{input\_records}: 228
  \item \texttt{post\_merge\_te\_nodes}: 225
  \item \texttt{homepage\_chart\_nodes}: 154
  \item \texttt{operation\_counts}: rename 27, merge 3
\end{itemize}

顶层 \texttt{counts} 已改为合并后的计数，即 \texttt{standard=138}, \texttt{A=36}, \texttt{B=35}, \texttt{C=16}。

\section{MySQL Expression 数据库}

\subsection{连接目标}

Expression MySQL 数据库由 \verb|api/expression_data.php| 读取配置。当前本地配置为：

\begin{itemize}
  \item host: \texttt{127.0.0.1}
  \item port: \texttt{3306}
  \item database: \texttt{tekg\_expression}
  \item charset: \texttt{utf8mb4}
\end{itemize}

\subsection{表计数}

\begin{longtable}{lr}
\toprule
Table & Rows \\
\midrule
\endhead
expression\_context\_catalog & 143 \\
expression\_dataset\_summary & 113604 \\
expression\_context\_stats & 5415124 \\
expression\_browse\_summary & 37868 \\
\bottomrule
\end{longtable}

\subsection{数据集}

\begin{longtable}{lllr}
\toprule
dataset\_key & dataset\_label & dataset\_order & Contexts \\
\midrule
\endhead
normal\_tissue & Normal Tissue & 1 & 32 \\
normal\_cell\_line & Normal Cell Line & 2 & 54 \\
cancer\_cell\_line & Cancer Cell Line & 3 & 57 \\
\bottomrule
\end{longtable}

Expression summary 当前覆盖 37868 个 TE 名称和 3 个 dataset。

\subsection{主要表用途}

\begin{longtable}{L{5.4cm}L{9.2cm}}
\toprule
Table & Purpose \\
\midrule
\endhead
expression\_context\_catalog & dataset/context catalog, context display order, and run count metadata. \\
expression\_dataset\_summary & per-TE per-dataset summary, including top contexts and aggregate max/mean/median metrics. \\
expression\_context\_stats & per-TE per-context statistics, including min, q1, median, q3, mean, max, std, and cv. \\
expression\_browse\_summary & one-row-per-TE browse table used by the Expression page. \\
\bottomrule
\end{longtable}

\subsection{文件来源}

Canonical expression asset root:

\begin{verbatim}
data/bulk_expression_web
\end{verbatim}

Important files:

\begin{itemize}
  \item \verb|data/bulk_expression_web/processed/te_expression_context_stats.tsv|
  \item \verb|data/bulk_expression_web/processed/te_expression_dataset_summary.tsv|
  \item \verb|data/bulk_expression_web/processed/te_expression_browse_summary.tsv|
  \item \verb|data/bulk_expression_web/processed/te_expression_context_catalog.tsv|
  \item \verb|data/bulk_expression_web/processed/expression_mysql_schema.sql|
\end{itemize}

Expression 构建脚本：

\begin{verbatim}
scripts/build/prepare_expression_assets.py
\end{verbatim}

\section{主要运行时 API}

不含智能问答时，核心数据库相关 API/入口为：

\begin{longtable}{lp{10cm}}
\toprule
File & Purpose \\
\midrule
\endhead
\path{api/runtime_config.php} & shared Neo4j runtime config loader. \\
\path{api/health.php} & health check, includes resolved Neo4j database name. \\
\path{api/taxonomy.php} & Neo4j-backed TE taxonomy API. \\
\path{api/taxonomy_lib.php} & shared PHP taxonomy query/helper library. \\
\path{api/graph.php} & graph search/query API, now includes taxonomy payload for TE nodes. \\
\path{api/te_metrics.php} & degree-style TE metric API. \\
\path{api/expression_data.php} & MySQL expression data access helper. \\
\bottomrule
\end{longtable}

\section{主要页面与数据库关系}

\begin{longtable}{L{4.8cm}L{9.8cm}}
\toprule
Page & Database/data dependency \\
\midrule
\endhead
\path{index.php} & Uses Neo4j taxonomy for homepage ring chart; falls back to \path{data/processed/tekg3_homepage_taxonomy.json} only if live query fails. \\
\path{preview.php} & Uses graph API/runtime graph data. \\
\path{search.php} & Uses Repbase, Dfam, JBrowse assets, and Neo4j-backed taxonomy helper. \\
\path{browse.php} & Uses Neo4j-backed taxonomy helper and local browse assets. \\
\path{genomic.php} and \path{jbrowse.php} & Use JBrowse assets under \path{data/JBrowse}. \\
\path{expression.php} and \path{expression_detail.php} & Use MySQL \texttt{tekg\_expression} and fallback TSV under \path{data/bulk_expression_web}. \\
\path{epigenetics.php} & Static/placeholder-oriented runtime page at present; future database dependency should be documented before adding data. \\
\bottomrule
\end{longtable}

\section{运行时数据资产}

\begin{longtable}{L{6.2cm}L{8.4cm}}
\toprule
Path & Role \\
\midrule
\endhead
\path{data/taxonomy/transposon_tree/tree_rmsk_repbase.txt} & TE taxonomy source input used by the standardization workflow. \\
\path{data/taxonomy/transposon_tree/tree_all.txt} & broader TE taxonomy source input. \\
\path{data/taxonomy/te_234/*} & TE taxonomy/reference material. \\
\path{data/processed/tekg3_taxonomy_standardization_report.json} & generated taxonomy standardization report. \\
\path{data/processed/tekg3_homepage_taxonomy.json} & homepage taxonomy fallback cache, not primary runtime source. \\
\path{data/processed/dfam/dfam_curated_catalog.json} & Dfam curated catalog used by search/detail flows. \\
\path{data/bulk_expression_web/*} & expression raw and processed assets. \\
\path{data/JBrowse/*} & genome browser and repeat locus assets. \\
\path{data/terminology/*} & terminology assets. \\
\bottomrule
\end{longtable}

\section{构建与检查脚本}

\begin{longtable}{L{6.2cm}L{8.4cm}}
\toprule
Script & Purpose \\
\midrule
\endhead
\path{scripts/normalize/build_tekg3_from_tekg21.py} & builds/normalizes \texttt{tekg3} from \texttt{tekg21} and writes taxonomy outputs. \\
\path{scripts/build/prepare_expression_assets.py} & builds expression TSV summaries and MySQL schema from \path{data/bulk_expression_web}. \\
\path{scripts/checks/check_taxonomy_runtime_consistency.py} & validates Neo4j-backed taxonomy runtime consistency. \\
\path{scripts/checks/check_runtime_db_config.py} & validates runtime DB config does not fall back to old DB names. \\
\path{scripts/checks/check_expression_paths.py} & validates expression path strategy uses \path{data/bulk_expression_web}. \\
\bottomrule
\end{longtable}

\section{Current Caveats}

\begin{itemize}
  \item \texttt{BIO\_RELATION} predicate vocabulary is not yet strongly normalized.
  \item Expression data is large and should stay summarized in Neo4j rather than fully duplicated into the graph.
  \item Some historical scripts still reference old \texttt{tekg2}, \texttt{tekg21}, or old taxonomy filenames; these should remain clearly legacy.
  \item Some Python bytecode files are currently tracked and should be removed from Git history/index in a cleanup pass.
  \item Large runtime data should receive an explicit manifest with path, role, source, size, and generating script.
\end{itemize}

\end{document}
